% Template LaTeX para Informes de Proyectos de Grado de Computación
% InCo, Facultad de Ingeniería, Universidad de la República
% 06/2023

\documentclass{../../prgrado}



%%%%%%%%%%%%%%%%%%%%%%%%%%%%%%%%%%%%%%%%%%%%%%%%%%%%%%%%%%%%%%%%%%%%%%%%%%%%%%%%%%%%%%%%%%%%%%%%
% Datos del Proyecto:
%%%%%%%%%%%%%%%%%%%%%%%%%%%%%%%%%%%%%%%%%%%%%%%%%%%%%%%%%%%%%%%%%%%%%%%%%%%%%%%%%%%%%%%%%%%%%%%%

% título del proyecto (debe escribirse en minúscula, con excepción de la
% letra inicial de la primera palabra y nombres propios)
\title{Modulo Generativo basado en \Gls{OpenFing}.}                   

% autores
\author{Diego Amorena, Guido Dinello}  

% fecha de la defensa
\date{\today}                                 

% supervisor
\supervisor{Luis Chiruzzo}                

% cosupervisor (comentar si no tiene)
\cosupervisor{Libertad Tansini}

% licencia creative commons del documento
% opciones: by, by-sa, by-nd, by-nc, by-nc-sa y by-nc-nd
% opción por defecto: by-nc-nd versión 4.0
\cclicense{by}{4.0}

%%%%%%%%%%%%%%%%%%%%%%%%%%%%%%%%%%%%%%%%%%%%%%%%%%%%%%%%%%%%%%%%%%%%%%%%%%%%%%%%%%%%%%%%%%%%%%%%


%links con colores 
\hypersetup{
  colorlinks   = true,
  urlcolor     = blue,
  linkcolor    = blue,
  citecolor    = red
}


\begin{document}

%%%%%%%%%%%%%%%%%%%%%%%%%%%%%%%%%%%%%%%%%%%%%%%%%%%%%%%%%%%%%%%%%%%%%%%%%%%%%%%%%%%%%%%%%%%%%%%%
% Parte inicial
%%%%%%%%%%%%%%%%%%%%%%%%%%%%%%%%%%%%%%%%%%%%%%%%%%%%%%%%%%%%%%%%%%%%%%%%%%%%%%%%%%%%%%%%%%%%%%%%

\frontmatter % numeración en romanos, capitulos sin numerar

% carátula
\maketitle


%%%%%%%%%%%%%%%%%%%%%%%%%%%%%%%%%%%%%%%%%%%%%%%%%%%%%%%%%%%%%%%%%%%%%%%%%%%%%%%%%%%%%%%%%%%%%%%%

% agradecimientos
\chapter*{Agradecimientos}

Agradecer, siempre es bueno agradecer.
%\textcolor<2>{orange}{Agradecimientos a Libertad Tansini y Luis Chiruzzo}



\chapter*{Resumen}
\begin{itemize}
    \item Una breve mención del estado del arte. Donde mencionaremos que las herramientas de inteligencia artificial se han vuelto esenciales en muchos ámbitos tanto profesionales como educativos.
    \item La falencias que tiene estas herramientas, que no siempre usan fuentes confiables y tienden a alucinaciones. 
    \item La necesidad, potencial y el valor de una herramienta, confiable y especifica, que apoye al estudiante de la FING en su aprendizaje durante su carrera educativa por la facultad.
    \item ¿Que debería incluir una herramienta como la mencionada?
    Mencionar que debería de:
    \begin{itemize}
        \item Debe incluir una base vectorial con los cursos de \Gls{OpenFing}.
        \item Debe incluir un generador de embeddings, tanto para prompts como para chunks/cursos.
        \item Debe incluir un generador una búsqueda de recuperación contextual.
        \item Debe incluir una generación de respuestas a partir de prompts.
    \end{itemize}

    \item ¿Como evaluaríamos si la herramienta cumple lo especificado?
    \begin{itemize}
        \item Debe responder coherentemente en base a los cursos de \gls{OpenFing}.
        \item Debe referenciar al contexto  utilizado. Para verificar la respuesta o profundizar en el tema. 
        \item Debe dejar en claro que no se encontró el contexto requerido cuando es el caso.
    \end{itemize}
\end{itemize}


\hfill \break
\keywords{A decidir}

%%%%%%%%%%%%%%%%%%%%%%%%%%%%%%%%%%%%%%%%%%%%%%%%%%%%%%%%%%%%%%%%%%%%%%%%%%%%%%%%%%%%%%%%%%%%%%%%

% índice
\tableofcontents
\newpage


%%%%%%%%%%%%%%%%%%%%%%%%%%%%%%%%%%%%%%%%%%%%%%%%%%%%%%%%%%%%%%%%%%%%%%%%%%%%%%%%%%%%%%%%%%%%%%%%
% Parte central
%%%%%%%%%%%%%%%%%%%%%%%%%%%%%%%%%%%%%%%%%%%%%%%%%%%%%%%%%%%%%%%%%%%%%%%%%%%%%%%%%%%%%%%%%%%%%%%%

\mainmatter %numeración en arábicos, numerar capítulos 

% introducción
\chapter{Introducción}

En la introducción se deberá:
\begin{itemize}
    \item  Explayar en el Estado del Arte.
    \item Explicar algunas en detalle algunas de las posibles falencias o problemas que se generar al utilizar modelos que no fueron entrenados con contenido educativo o que no fueron entrenados con contenido local/regional. 
    Ejemplo de problemas: nombres, enfoque o Convenciones, distintos que nosotros(Udelar o Uruguay) podamos usar comparado a un ámbito internacional. 
    Exceso de confianza por los modelos a sus respuesta. 
    Ejemplo: Un casó de valor puede ser que el Metodo Simplex(Introducción a la Investigación Operativa) se da de manera distinta.
    \item Motivación del proyecto. Mencionar potencial y valor que le añadiría a los estudiantes. 
    \item Definición de una herramienta RAG
    \item Mencionar los dos Proyectos previos que realizaron en WEBIR.
    \item Se deberá detallar el scope del proyecto, el alcance que nos propusimos. Se deberá detallar cada uno de los parámetros de entrada y de los parámetros que se investigaron y evaluaron para llegar a un mejor resultado.
    \item Se deberá detallar en base a que evaluamos un modelo y en que nos basaremos para fundamentar que un resultado sea mejor o peor que otro.
    \item  Mencionar el procesos que se llevó a cabo.
    \item Mencionar si el desarrollo de dicha herramienta fue exitoso.
\end{itemize}
%%%%%%%%%%%%%%%%%%%%%%%%%%%%%%%%%%%%%%%%%%%%%%%%%%%%%%%%%%%%%%%%%%%%%%%%%%%%%%%%%%%%%%%%%%%%%%%%



% antecedentes
\chapter{Revisión de antecedentes}

Acá se pueden revisar casos uruguayos de herramientas RAG. 
En caso de no encontrar ningún caso uruguayo, se podrán detallar casos educativos similares.


%%%%%%%%%%%%%%%%%%%%%%%%%%%%%%%%%%%%%%%%%%%%%%%%%%%%%%%%%%%%%%%%%%%%%%%%%%%%%%%%%%%%%%%%%%%%%%%%

\chapter{Marco Teórico}

Aquí se explicaran cada uno de los conceptos utilizados en el documento.

\section{Conceptos Generales}
Conceptos necesarios generales que se precisan para entender el proyecto.
\subsection{Tokenizacion}
\subsection{Modelo de Lenguaje}
\subsection{Red Neuronal}
\subsection{Red Neuronal de Aprendizaje Profundo}
\subsection{Grandes Modelos de Lenguage}
\
\subsection{Promt}
\subsection{Embbeding}
\subsection{Chunks}

\subsection{Retrieval-Augmented Generation(RAG)}

\section{Estadísticas de Evaluación}
Acá se definirá cada una de las estadísticas que se utilizaran posteriormente en Evaluación/Experimentación.
\subsection{Word Error Rate}
\subsection{Word Error Rate}

\section{Antecedentes Relevantes}

\subsubsection{OpenFing}


%%%%%%%%%%%%%%%%%%%%%%%%%%%%%%%%%%%%%%%%%%%%%%%%%%%%%%%%%%%%%%%%%%%%%%%%%%%%%%%%%%%%%%%%%%%%%%%%


\chapter{Desarrollo}
Aquí se mencionará procesos que no entran dentro del producto final pero fueron necesarios para el desarrollo de la misma. 

\section{Scraping de OpenFing}
Descripción de lo que involucra este proceso.

\subsection{Problemas encontrados}
Acá se puede incluir los errores que se encontraron durante el scraping del proceso. 

Por ejemplo:

La mayoría de videos siguen esta estructura, pero no todos.\\
\url{https://open.fing.edu.uy/media/Acronimo del curso/Acronimo del curso_Numero_de_Clase.mp4}

El proceso de scraping se hizo ligeramente más largo para evitar este problema de buscar video en urls equivocados.

\subsection{Mejoras Implementadas}

Aquí se detallaran las mejoras que recibió el proceso a partir de los problemas. 
Por ejemplo: 

Se empezó a ingresar a cada una de las paginas de los videos, y no solo a la pagina del curso.
\subsection{Mejoras Posibles}
Aquí se detallaran las mejoras pendientes.
Por ejemplo: 
Cambiar el método de la descarga dependiendo de si el host es youtube o \gls{openfing}.

\section{Transcripción de clases}

Descripción de lo que involucra este proceso.
\subsection{Problemas encontrados}
El modelos de whisper python es:
    Muy lento.
    Repite el texto anterior cuando encuentra silencio. 
    Problemas de implementación en Docker cuando no se tiene disponible una GPU (pythorch).
\subsection{Mejoras Implementadas}
Cambio de modelos a whisper.cpp y faster-whisper. Si se transforma el video a otros formatos se puede eliminar el video y transcribir mar rápido. 
\subsection{Mejoras Posibles}


%%%%%%%%%%%%%%%%%%%%%%%%%%%%%%%%%%%%%%%%%%%%%%%%%%%%%%%%%%%%%%%%%%%%%%%%%%%%%%%%%%%%%%%%%%%%%%%%


\chapter{Tecnologías}
Aquí se detallaran las tecnologías que se decidieron utilizar como también aquellas que fueron descartadas.
\section{Tecnologías Utilizadas}

Se tendrá un listado de las tecnologías seleccionadas. Sumado con una breve descripción de que son y para que se utilizaron dentro del proyecto.
\subsection{Lenguajes}
\subsection{Método de Transcripción}
\subsection{Generador de embbedings}
\subsection{Base de datos vectorial}
\subsection{Generador de respuestas}
\subsection{Frontend}
\section{Tecnologías Descartadas}

Se tendrá un listado de las tecnologías utilizadas pero descartadas durante el proceso de desarrollo o previo al mismo con el motivo por el cual se terminaron descartando.  

\subsection{Lenguajes}
\subsection{Método de Transcripción}
\subsection{Base de datos vectorial}
\subsection{Generador de embbedings}
\subsection{Frontend}


\chapter{Arquitectura}
Aquí se tendrá un diagrama uml de la arquitectura final del proyecto.
\section{Diseño}
Aqui se digramará la arquitectura y se incluira una  breve descripcion en lo necesario. 
\section{Decisiones de Arquitectura}
Aquí se explicaran algunas de las decisiones que se tomaron acerca de la arquitectura del proyecto. 
\subsection{¿Arquitecturas Descartadas?Opcional}
Aquí se puede explicar por que se descartaron ciertas arquitecturas que fueron candidatas.


\chapter{Funcionalidades realizadas}
En esta sección se discutirá cada una de las funcionalidades realizadas. 
\section{Scraping de OpenFing}
En esta sección se discutirá las funcionalidades que no llegaron a hacer que entran dentro del scope
\subsection{Funcionalidades no culminadas}
\section{Descargado de Clases}
\section{Transcription}
\section{Generator}
%%%%%%%%%%%%%%%%%%%%%%%%%%%%%%%%%%%%%%%%%%%%%%%%%%%%%%%%%%%%%%%%%%%%%%%%%%%%%%%%%%%%%%%%%%%%%%%%

% experimentación


\chapter{Experimentación}
Aquí se detalla las experimentaciones que se realizaron. 
Involucra definición de cada uno de los parámetros de entrada y \gls{contexto} de donde se utilizaría en nuestro proyecto. 

\section{Modelos de Transcripción}
Benchmarks para los modelos de transcripción.
\subsection{Método de Evaluación}
Para esta sección se transcribieron ciertas clases y en base a dichas transcripciones se evaluaron los modelos. Ademas se comparo el tiempo de ejecución de la transcripción de distintos modelos y distintas librerías. 

\subsection{Modelos}
\subsubsection{Resultados Modelo Tiny}
\subsubsection{Resultados Modelo Base}
\subsubsection{Resultados Modelo Small}
\subsubsection{Resultados Modelo Medium}
\subsubsection{Resultados Modelo Large}
\subsection{OpenAI-Whiper}
\subsection{Whisper.cpp}
\subsection{faster-whisper}
\subsection{Visualización}
Gráfica de los modelos y librerias que compara las estadisticas motradas de manera visual.

\subsection{Conclusiones}
Evaluación de resultados de las metricas y del tiempo. 



\section{Recuperación de Búsqueda}
\subsection{Resultados}
Anexo a ejemplo.

\subsection{Conclusiones}

\section{Modelos de Generación}
Anexo a ejemplo.


\subsection{Conclusiones}

\subsection{Métricas}
\subsection{Resultados}
Análisis de las respuestas
Cuantitativo y Cualitativo.


%%%%%%%%%%%%%%%%%%%%%%%%%%%%%%%%%%%%%%%%%%%%%%%%%%%%%%%%%%%%%%%%%%%%%%%%%%%%%%%%%%%%%%%%%%%%%%%%

%conclusiones
\chapter{Conclusiones y Trabajo Futuro}

Se mencionara como concluyo el proyecto.
Se mencionara como mejorar el proyecto.
Se mencionara como mejorar el proceso de desarrollo.
Se mencionara que hubiesemos cambiado si empezásemos de 0.


%%%%%%%%%%%%%%%%%%%%%%%%%%%%%%%%%%%%%%%%%%%%%%%%%%%%%%%%%%%%%%%%%%%%%%%%%%%%%%%%%%%%%%%%%%%%%%%%
% Parte final
%%%%%%%%%%%%%%%%%%%%%%%%%%%%%%%%%%%%%%%%%%%%%%%%%%%%%%%%%%%%%%%%%%%%%%%%%%%%%%%%%%%%%%%%%%%%%%%%

{ % inicio backmatter

\backmatter % no numerar capítulos 

% referencias bibliográficas

\newpage
\bibliographystyle{apacite}
\bibliography{referencias}

} % fin backmatter

%%%%%%%%%%%%%%%%%%%%%%%%%%%%%%%%%%%%%%%%%%%%%%%%%%%%%%%%%%%%%%%%%%%%%%%%%%%%%%%%%%%%%%%%%%%%%%%%

% anexos
\begin{appendix}
\chapter{Organizacion}
\section{Recursos Utilizados}
\subsection{Repositorio}
\url{https://gitlab.fing.edu.uy/diego.amorena/tesis}
\end{appendix}
\begin{appendix}

\chapter{Anexos de herramienta}
\section{Manual de Usuario}

\section{Métodos expuestos}

\section{Testing Realizados}
\section{¿Known Bugs?}

\subsection{Tablero de Tareas}
\url{https://trello.com/invite/b/67d23719de5ec54f6976dcf2/ATTI115d1f2970fa4f083c6468d14b435b5e71CFEFF9/backlog}
\end{appendix}
\begin{appendix}
\chapter{Resultados}
Aquí se incluyen resultados varios
\end{appendix}



%%%%%%%%%%%%%%%%%%%%%%%%%%%%%%%%%%%%%%%%%%%%%%%%%%%%%%%%%%%%%%%%%%%%%%%%%%%%%%%%%%%%%%%%%%%%%%%%

\end{document}
